\problemname{Leaky Pipe}
Bob the plumber works in a massive hydroelectric power plant, where massive pipes direct the flow of water through the plant.
One pipe would flow into a manifold, which would then disperse the water into other pipes.
Eventually all pipes would exhaust back into the river.
This increased the water flow throughout the entire plant, and allowed for redirecting water during maintainace.

However, today, disaster had struck, one of the pipes had begun to leak and it was important for it to be replaced as soon as possible.
Bob found the piping and instrumentation diagram and quickly located the leaky pipe.
As he looked at the pipe, he found the size to be a bit awkward, as it seemed oodly large for the pipes around it.

Now Bob wants to use this oppurtunity to decrease the size of the pipe or completely remove it, if at all possible.

\begin{Input}
The first line of input contains 6 integers:

\begin{itemize}
 \item $1 \le p \le 1000$: the number of pipes
 \item $\left\lceil \frac{1+\sqrt{1+4p}}{2} \right\rceil \le m \le p+1$: the number of manifolds
 \item $s$ the source for the pipe system
 \item $t$: the river which all pipes exhaust into($t \ne s$)
 \item $m_s$: the manifold of origin for broken pipe, which Bob wants to minimize in size ($m_s \ne t$)
 \item $m_t$: the manifold of termination for the broken pipe, which Bob wants to ($m_t \ne m_s$)
\end{itemize}
What follows is $p$ lines containing 3 integers $v$, $u$, $v \ne u$, and $1 \le c \le 10^4$, indicating a pipe from manifold $v$ to manifold $u$ with a capacity of $c$.
\end{Input}

\begin{Output}
The output should contain exactly one integer, the minimum size of the pipe from $m_s$ to $m_t$, which does not impact the maximum water flow.
\end{Output}

